% Separate post-CP19 candidate; not a symmetry or nonvanishing assertion.
\subsection{The variation as a derivative of the trace-lift matrix}

\begin{lemma}[A typed derivative identity]
\label{lem:DE-A3-positive-cup-derivative-candidate}
\status{Candidate; complete algebraic proof written, independent audit pending.}
Keep the sharp-source and etale-double-cover hypotheses of
Lemma~\ref{lem:DE-A3-positive-index-double-corank-candidate}.
Write the source section parameters as squares, and denote their
square roots by $u$. Use compatible bases on the relative Frobenius
twists. There is a rational $2\times2$ matrix $A(u)$ on the sharp
source open, obtained by the actual etale trace lifts and projection
along the Cartier image, with the following property: the root-removed
trace-direction variation is exactly
\[
                         D^{\mathrm{root}}_v=dA_u(v).
                                                        \tag{DE-CD1}
\]
Here the actual coherent first-variation matrix has the corresponding
Frobenius-squared coordinates and the actual section direction is the
square of $v$. In particular, trace-zero of every variation is equivalent
to $d\operatorname{tr}A=0$. Over a perfect ground field this says that
$\operatorname{tr}A$ is a square in the source parameter function field.
Neither that square condition nor determinant nonvanishing is asserted.
\end{lemma}
\begin{proof}
We retain all twists in the construction. Let $U$ be the vector space
of square-root section coordinates and write $W=U^{(2)}$ for the
ordinary section-coordinate space. Root-removed Cartier multiplication
is a symmetric alternating bilinear map
\[
             B:U\otimes U\longrightarrow W.
\]
The alternating assertion is the mixed-volume identity already proved.
Both $U$ and $W$ have compatible decompositions into the two toral
weights. The map $B$ is zero between different weights. Put
$M(u)v=B(u,v)$, write $u_i$ for the weight-$i$ part of $u$, and set
\[
 r(u)=(u_0,u_1),\qquad R(u)=(u_0^{[2]},u_1^{[2]}).
\]
Sharpness gives the direct sum $W=\im M(u)\oplus\im R(u)$.
Let $P(u):W\to k^2$ be its projection in the displayed $R$ basis.
Thus
\[
                       PM=0,\qquad PR=1.                \tag{DE-CD2}
\]
All these matrices are rational functions on the sharp open.

Choose a toral-compatible linear trace splitting
$q:U\to\widetilde U$ of the etale cover; use its Frobenius twist
$q^{(2)}:W\to\widetilde W$ on ordinary coordinates. The deck
involution sends $qv$ to $qv+\pi^*v$, and similarly on $W$.
Cartier and trace commute. There are therefore bilinear maps
$H,J,K:U\otimes U\to W$ such that
\[
 \begin{aligned}
 \widetilde B(\pi^*u,qv)&=\pi^*H(u,v)+q^{(2)}B(u,v),\\
 \widetilde B(qu,qv)&=\pi^*J(u,v)+q^{(2)}K(u,v).
 \end{aligned}
                                                        \tag{DE-CD3}
\]
Equivariance under the deck involution, followed by expansion of
both sides, gives the exact identity
\[
                  K(u,v)=H(u,v)+H(v,u)+B(u,v).            \tag{DE-CD4}
\]
The maps $H$ and $B$ vanish on different toral weights.

Define the trace-lift coefficient matrix by
\[
                          A(u)=P(u)H(u)r(u),
                                                        \tag{DE-CD5}
\]
where $H(u)v=H(u,v)$. Choose, locally on the sharp open, columns
$w_i\in U$ with
\[
                 H(u,u_i)+M(u)w_i=R(u)A(u)_i.            \tag{DE-CD6}
\]
Such columns exist by (DE-CD2). The square-root coordinates
$q u_i+\pi^*w_i$ are precisely the corrected trace lifts used in
Lemma~\ref{lem:DE-A3-positive-trace-radical}: their ordinary
coordinates have trace $u_i^{[2]}$, and their Cartier images belong
to $\pi^*\im R$. The root-removed variation in trace direction $v$
therefore has column
\[
                  P\bigl(K(v,u_i)+B(v,w_i)\bigr).
                                                        \tag{DE-CD7}
\]
This formula is the earlier trace-dual kernel--cokernel formula,
with its semilinear types retained.

Differentiate (DE-CD2) with respect to the root parameters $u$.
In characteristic two $dR=0$, so
\[
                  (dP)R=0,\qquad(dP)M=P\,dM.
\]
Using (DE-CD6) gives
\[
 \begin{aligned}
 dA_u(v)_i
 &=P B(v,w_i)+P H(v,u_i)+P H(u,v_i).
 \end{aligned}
                                                        \tag{DE-CD8}
\]
Toral orthogonality gives $H(u,v_i)=H(u_i,v)$ and
$B(v,u_i)=B(u,v_i)\in\im M(u)$, the last term being killed by $P$.
Substitution of (DE-CD4) into (DE-CD7) now proves (DE-CD1), column by
column. The identities also show that choices of the correction
columns do not change the induced variation.

Taking traces yields the asserted scalar criterion. The sharp open
is an open subset of the section affine space, whose function field
is a rational function field over a perfect field. In such a field
the kernel of the absolute differential is exactly the subfield of
squares: expand in the basis of square-free monomials over the square
subfield and differentiate. This proves the final equivalence.
\end{proof}

\begin{remark}[Diagnostic evidence and remaining scope]
\label{rem:DE-A3-positive-cup-derivative-scope-candidate}
\status{Candidate scope statement; uniform symmetry remains unproved.}
The retained programs
\path{agents/obstructions_crystals/code/post_CP19_canonical_cup_variation_probe.py},
\path{agents/obstructions_crystals/code/post_CP19_positive_local_index_cup_probe.py}
and
\path{agents/obstructions_crystals/code/post_CP19_GF4_canonical_variation_probe.py}
use the actual Cartier-image complement, trace lifts and Frobenius
linearization. Their author runtimes were respectively
$0.022789923001255374$, $0.10455728400120279$ and
$0.16540274800354382$ seconds. The first two retain twenty examples
each, of source types $(6,3)$ and $(6,6)$, and the last retains twelve
examples over $\mathbf F_4$. All sampled canonical variation matrices
have trace zero. The second family has positive local polarization
index. These finite computations neither establish the square
condition on arbitrary mixed surfaces nor exclude a common
annihilator. The lemma itself is a linear-algebra identity, not an
inference from those samples.
A separate direct evaluation over $\mathbf F_4[\epsilon]/(\epsilon^2)$
checks (DE-CD1) in all $216$ coordinate directions of the twelve retained
sources. It evaluates the actual Cartier products and inverts the
projection matrix over dual numbers, without using the differentiated
projection formula. The retained program
\path{agents/obstructions_crystals/code/post_CP19_GF4_cup_derivative_check.py}
passed in $1.857703956993646$ seconds.
\end{remark}
